% Pursuing Stacks (cotes 134-2 to 134-8): archive cote 134-6, site page 123 (PDF page 124), typescript page 441.
% typescript. Draft compared: ps5.tex:7948-8009.
% AI-transcribed by Claude Opus 5.5 (claude-opus-5-5, adaptive thinking, effort high), Message Batches API, batch ps3 (msgbatch_01XdFmuVVBAcgc7EFxdtCqsj), 2026-11-12; not yet checked by eye.
\scanpage{134-6/123}
\mspage{441}

involving a group (or better still, a groupoid) and an operation of \xout{the}
it on a structure of type (\unsure{$\ast$}) (embodying a universal covering...).

Here I am getting though into thin air again, and I do'nt expect
I'll ponder much more in this direction and see what comes out. The
striking fact however \xout{that} here is that quite unexpectedly, \unsure{We|we} got fur-%
ther hold and food for this thin-air intuition (which came up first in
relation to De Rham structures with devided powers), that there may be
a reasonable (and essentially just one such) notion of a ``homotopy type
over the ground ring k'' for any commutative ring k, reducing for $k=\underline{Z}$
to usual homotopy types, and giving rise to base change functors
\[ \mathrm{Hot}(k) \longrightarrow \mathrm{Hot}(k') \]
for any ring homomorphism $k\to k'$. And I wonder whether this might not
come out in some very simplistic way, in the general spirit of our
``modelizing story'', without having to work out in full a description
of homotopy types by such sophisticated models as De Rham complexes
with devided powers, or coalgebra structures in (Cat), and \inserted{looking up} maybe the
relation\hand{s} between these. (How by all means hope to recover a De Rham
structure from a stupid structure (\unsure{$\ast$}) ???)

2) The condition d) in corollary 1 is needed in order to ensure
\xout{\unsure{2)xx}} that a derived functor $\mathrm{RHom}_k(L\hand{$_\bullet$},F\hand{$^\bullet$})$ may be defined using \uline{projective} reso-%
lutions of L. , whereas conditions \xout{a)}b)\hand{,}c) ensure that a functor
$\mathrm{RHom}_k$ may be defined using \xout{proje} \uline{injective} resolutions of $K^\bullet$ . It is
a well known standard fact of homological algebra that in case both
methods work (namely here, when all four assumptions are satisfied)
that the two methods yield the same result, which may equally be
described by resolving simultaneously the two arguments. (NB condition
a) is needed anyhow for $\mathrm{Hom}_k$ to be defined and for $P^{\&}_{\hand{$\mathcal{M}$}}$ being an abelian
category, which allows to define $D^+(P^{\&}_{\hand{$\mathcal{M}$}})$.) Our preference goes to
the first method, which in case \xout{k=k} \hand{$P=\mathrm{Add}_k(A)$} and $L = k_A$ , conduces to
computation of cohomology $R\Gamma(F^\bullet)$ in terms of a ``cointegrator'' $L_\bullet^A$ on
A . However, when it comes to introducing the variant $R\underline{\mathrm{Hom}}_k$ , this
method breaks down, as we saw, it is the other one which works. We
thus get a satisfactory formalism of $\mathrm{RHom}_k$ and $R\underline{\mathrm{Hom}}_k$ (including formula
(134) relating them via $RH^\bullet$ ) using only assumptions a)b)c).

3) If we want to extend the $R\underline{\mathrm{Hom}}_k$ formalism to the set-up when
the data A is replaced by a k-additive category \inserted{P} endowed with a diagonal
map as in remark 1), the proof on page 437 shows that what is needed

\margin{123}
